% !TEX program = xelatex
% Шаблон научной статьи НИТУ «МИСиС»
% Заполните блоки ниже и скомпилируйте: xelatex article.tex

\documentclass{misisarticle}

% --- Классификационные коды (оставьте пустым {}, если не требуются) ---
\grnti{28.23.15, 47.01.09}
\asjc{1707, 2207}
\oecd{1020, 1070}

% --- Название ---
\title{Использование дополнительного зеркального элемента на наземном роботе для повышения точности лидарной одометрии}

% --- Авторы: ---
\misisauthor{М.Г. Шитова}{кафедра инженерной кибернетики НИТУ «МИСиС»}{Москва, Россия}{m2302964@edu.misis.ru}
\misisauthor{М.Р. Вишневский}{кафедра инженерной кибернетики НИТУ «МИСиС»}{Москва, Россия}{supervisor@misis.ru}

% --- Аннотация и ключевые слова ---
\abstract{%
  Современные наземные робототехнические системы требуют надёжной пространственной ориентации в реальном времени, и ключевой технологией здесь выступают лидарная одометрия и SLAM. Качество их работы критически зависит от плотности облака точек, однако возможности механических лидаров ограничены числом лазерных каналов и частотой вращения, а требования к массогабаритным характеристикам платформ ограничивают замену сенсора. Перспективным решением является размещение дополнительного пассивного зеркального элемента в поле зрения лидара: отражённый сигнал создаёт виртуальный канал, повышающий плотность и расширяющий поле обзора сканируемого пространства. Однако вопрос о том, как такое искусственное увеличение плотности влияет на точность и робастность лидарной одометрии наземного робота, остаётся открытым: с одной стороны, дополнительные точки увеличивают число геометрических соответствий, с другой~--- вносят искажения, связанные с ошибками позиционирования зеркала. Целью работы является исследование влияния дополнительного зеркального элемента на точность лидарной одометрии и оценка этого эффекта на алгоритме FAST-LIO. В работе изучены принципы построения лидарной одометрии и SLAM, развёрнуто рабочее окружение и проведены эксперименты на четырёх эталонных наборах данных HKU в indoor- и outdoor-сценах. Проведено исследование влияния синтетического зеркалирования облака точек на точность одометрии с использованием метрики start-end position error. Полученные количественные оценки могут быть использованы при проектировании сенсорных систем наземных роботов с пассивными оптическими элементами.
}

\keywords{%
  лидар, зеркальное сканирование, облако точек, FAST-LIO, лидарная одометрия, SLAM
}

\begin{document}

\maketitle

% ============================================================
\section{Введение}

Современные наземные робототехнические системы~--- от беспилотных автомобилей до складских роботов и инспекционных дронов~--- требуют надёжной пространственной ориентации в реальном времени. Ключевой технологией, обеспечивающей решение этой задачи, является лидарная одометрия и SLAM: методы, позволяющие устройству определять собственное положение и одновременно строить карту окружающей среды, опираясь на данные бортовых сенсоров~\cite{thrun2005,xu2021}. Основным источником информации в таких системах выступает лидар (LiDAR, Light Detection and Ranging)~--- устройство, формирующее трёхмерное облако точек сцены посредством лазерного сканирования.

Качество работы алгоритмов лидарной одометрии и SLAM критически зависит от характеристик входного облака точек, в числе которых одним из ключевых является плотность точек~--- количество измерений на единицу площади сканируемой поверхности~\cite{beraldin2010}. В условиях ограниченной разрешающей способности механических лидаров~--- наиболее распространённого типа сенсоров в мобильной робототехнике~--- и жёстких требований к массогабаритным характеристикам платформ возникает задача повышения информативности облака точек без замены основного сенсора~\cite{wang2020,yoo2018}.

Одним из перспективных технических решений является применение зеркального сканирования~--- размещение в поле зрения лидара пассивного отражающего элемента (плоского зеркала), создающего дополнительный виртуальный канал и тем самым увеличивающего плотность и расширяющего поле обзора сканируемого пространства~\cite{vishnevsky2026}. В работе~\cite{vishnevsky2026} представлен метод повышения плотности LiDAR-сканирования с использованием пассивной зеркальной оптики и алгоритма автоматической калибровки GPC (Ground Plane Consistency). Математическая модель отражения базируется на операторе Хаусхолдера~--- матрице $H \in \mathbb{R}^{4\times4}$, параметризованной единичной нормалью $n$ к плоскости зеркала и расстоянием $d$ от оптического центра сенсора. Экспериментально показано, что зеркальный канал обеспечивает прирост заполненности до 50\,\% (в симуляции) и до 40\,\% (в натурных условиях) при сохранении низкой стоимости аппаратной реализации~\cite{vishnevsky2026,lidar3d2025}.

Однако вопрос о том, каким образом искусственное увеличение плотности облака точек методом зеркального сканирования влияет на точность и робастность алгоритмов лидарной одометрии и SLAM в динамическом режиме, остаётся открытым. С одной стороны, дополнительная плотность потенциально способна улучшить регистрацию последовательных сканов за счёт большего числа геометрических соответствий. С другой~--- отражённые точки вносят геометрические искажения, связанные с ошибками позиционирования зеркала. Актуальность настоящего исследования определяется необходимостью систематического анализа влияния плотности зеркального сканирования на точность оценивания траектории алгоритмами лидарно-инерциальной одометрии.

Целью работы является исследование влияния дополнительного зеркального элемента на точность лидарной одометрии наземного робота и количественная оценка этого эффекта на алгоритме FAST-LIO. Для достижения поставленной цели решаются следующие задачи: изучить принципы построения алгоритмов лидарной одометрии и SLAM и роль плотности облака точек в обеспечении точности оценивания траектории; рассмотреть архитектуру и математическую основу алгоритма FAST-LIO; развернуть рабочее окружение и воспроизвести алгоритм на открытых датасетах HKU; провести предварительное исследование влияния синтетического зеркалирования облака точек на точность одометрии с использованием метрики start-end position error; сформулировать направления дальнейших исследований.

% ============================================================
\section{Лидарные сенсоры и облака точек}

Лидар представляет собой технологию оптического дальнометрирования, основанную на измерении времени пролёта (ToF, Time of Flight) лазерного импульса от источника до объекта и обратно к приёмнику. Зная скорость распространения света $c$, расстояние до объекта $R$ определяется по формуле $R = c\Delta t/2$, где $\Delta t$~--- время между излучением импульса и регистрацией отражённого сигнала. Современные лидары излучают сотни тысяч импульсов в секунду, формируя трёхмерное облако точек~--- множество координат $p_i=(x_i,y_i,z_i)$ в пространстве, описывающее геометрию окружающих объектов~\cite{beraldin2010}.

По принципу сканирования лидары подразделяются на два основных типа: механические вращающиеся и твердотельные. Механические лидары (например, Velodyne VLP-16, HDL-64E) используют вращающийся блок лазерных излучателей и приёмников, обеспечивая панорамный обзор $360^\circ$ в горизонтальной плоскости. Однако плотность облака точек таких сенсоров ограничена фиксированным числом лазерных каналов (16, 32, 64 или 128) и частотой вращения. Вертикальное разрешение механического лидара определяется угловым шагом между соседними каналами (типично $2^\circ$ для VLP-16), что приводит к существенной неравномерности покрытия: объекты на большом расстоянии сканируются с низкой плотностью~\cite{vishnevsky2026,beraldin2010,yoo2018}.

Лидары Livox (AVIA, MID-40, MID-70, Horizon) используют иной принцип~--- неповторяющийся паттерн сканирования (non-repetitive scanning pattern). Лазерный луч отклоняется с помощью призменной системы, описывая на поверхности объектов розеточную траекторию, которая не повторяется от кадра к кадру. С течением времени поле зрения покрывается всё более плотно: за 0,1~с покрытие FOV составляет порядка 60\,\%, за 0,5~с~--- около 90\,\%, за 1~с~--- порядка 99\,\%~\cite{xu2021,xu2022}. Данная особенность делает лидары Livox особенно привлекательными для задач SLAM, где накопление плотной карты происходит естественным образом по мере движения сенсора.

Облако точек (point cloud) представляет собой неупорядоченный набор точек в трёхмерном пространстве. Каждая точка $p_i$ характеризуется как минимум координатами $(x,y,z)$; в зависимости от сенсора могут быть доступны дополнительные атрибуты: интенсивность отражённого сигнала $I_i$, временная метка $t_i$, номер лазерного канала. В экосистеме ROS облака точек передаются через топики (topics) типа \texttt{sensor\_msgs/PointCloud2}. Для лидаров Livox используется собственный формат сообщений \texttt{livox\_ros\_driver/CustomMsg}, обеспечивающий более эффективную сериализацию данных специфического формата сканирования~\cite{xu2021,xu2022}.

Плотность облака точек является ключевой характеристикой, определяющей качество downstream-задач: регистрации сканов, детекции объектов, сегментации и построения карты. Одним из перспективных методов увеличения плотности без замены сенсора является зеркальное сканирование~--- размещение пассивного отражающего элемента в поле зрения лидара. В работе~\cite{vishnevsky2026} предложен подход, при котором плоское зеркало создаёт виртуальный канал, а алгоритм GPC автоматически калибрует параметры зеркала, обеспечивая точное совмещение прямого и отражённого каналов. Экспериментально показан прирост заполненности до 50\,\% в симуляции и до 40\,\% в натурных условиях~\cite{vishnevsky2026,lidar3d2025}. Близкие идеи применяются и для повышения покрытия при сканировании поверхностей~\cite{flatness2021}.

% ============================================================
\section{Лидарная одометрия и SLAM}

Задача лидарной одометрии (LiDAR Odometry, LO) состоит в оценке движения сенсора~--- его положения и ориентации в пространстве~--- на основе последовательности облаков точек, поступающих с лидара. Формально: дана последовательность сканов $\{S_1,S_2,\dots,S_k\}$, требуется восстановить траекторию $T=\{T_1,T_2,\dots,T_k\}$, где $T_i\in SE(3)$~--- преобразование из системы координат сенсора в момент $i$ в глобальную (или начальную) систему координат~\cite{thrun2005,zhang2014}. В отличие от одометрии на основе колёсных энкодеров, лидарная одометрия не накапливает ошибку из-за проскальзывания колёс и не зависит от рельефа поверхности.

Базовым методом регистрации (alignment) двух последовательных облаков точек является алгоритм ICP (Iterative Closest Point) и его многочисленные производные~\cite{besl1992,segal2009}. ICP итеративно решает задачу минимизации суммы квадратов расстояний между соответствующими точками двух сканов: на каждой итерации для точек исходного облака находятся ближайшие точки в целевом облаке, после чего вычисляется оптимальное жёсткое преобразование. Обобщённый ICP (Generalized ICP, GICP) расширяет базовый алгоритм учётом локальной геометрии поверхности (нормалей или ковариационных матриц) в окрестности каждой точки, что повышает точность регистрации в структурированных сценах~\cite{segal2009}.

SLAM (Simultaneous Localization and Mapping) расширяет задачу одометрии: помимо оценки текущего положения, алгоритм одновременно строит и уточняет карту окружающей среды. В отличие от чистой одометрии, SLAM явно поддерживает глобальную согласованность траектории: при обнаружении петли (loop closure)~--- возвращения сенсора в ранее посещённую область~--- выполняется коррекция накопленного дрейфа и обновление карты~\cite{thrun2005,zhang2014}. Современные лидарные SLAM-системы интегрируют данные инерциального измерительного модуля (IMU), что позволяет компенсировать быстрые движения и повысить робастность в сценах с недостатком геометрических признаков.

К числу наиболее известных алгоритмов лидарной одометрии и SLAM относятся:
\begin{itemize}
  \item \textbf{LOAM} (Lidar Odometry and Mapping, J.~Zhang, S.~Singh, 2014~\cite{zhang2014})~--- классический алгоритм, разделяющий задачу на высокочастотную одометрию (10~Гц) и низкочастотное построение карты (1~Гц). Основан на извлечении признаков (угловые и планарные точки) и минимизации расстояний «точка--линия» и «точка--плоскость».
  \item \textbf{LeGO-LOAM} (T.~Shan, B.~Englot, 2018~\cite{shan2018})~--- легковесная оптимизированная версия LOAM для наземных роботов. Включает сегментацию грунта, что повышает точность в структурированных outdoor-сценах.
  \item \textbf{LIO-SAM} (T.~Shan et al., 2020~\cite{shan2020})~--- тесно-связанная лидарно-инерциальная SLAM-система на основе факторных графов. Поддерживает замыкание петель, GPS-факторы и мультисенсорную интеграцию.
\end{itemize}

% ============================================================
\section{Алгоритм FAST-LIO}

FAST-LIO (Fast LiDAR-Inertial Odometry)~--- алгоритм тесно-связанной (tightly-coupled) лидарно-инерциальной одометрии, разработанный в Гонконгском университете (HKU) группой MARS Lab под руководством проф. Ф.~Чжана~\cite{xu2021}. Алгоритм выполняет совместную оптимизацию лидарных и инерциальных измерений в едином фильтре, обеспечивая высокую точность и вычислительную эффективность. Версия FAST-LIO2~\cite{xu2022} дополнительно использует прямую регистрацию точек (без извлечения признаков), что делает алгоритм инвариантным к паттерну сканирования сенсора.

\subsection{Архитектура и математическая основа}

FAST-LIO базируется на итерационном расширенном фильтре Калмана (Iterated Extended Kalman Filter, IEKF). Вектор состояния $x$ включает положение $p\in\mathbb{R}^3$, ориентацию (кватернион $q\in SO(3)$), скорость $v\in\mathbb{R}^3$, а также смещения (bias) акселерометра $b_a$ и гироскопа $b_g$ IMU. Измерения IMU используются на этапе предсказания (prediction) для распространения состояния между моментами поступления лидарных сканов с высокой частотой (200~Гц). На этапе коррекции (update) выполняется итерационное уточнение оценки состояния путём минимизации расстояний «точка--плоскость» между точками текущего скана и точками карты, накопленной в ikd-Tree~\cite{xu2021}.

Ключевой особенностью FAST-LIO, определяющей его вычислительную эффективность, является использование инкрементального kd-дерева (ikd-Tree)~--- динамической пространственной структуры данных для хранения и поиска точек карты~\cite{xu2021}. В отличие от статического kd-дерева, ikd-Tree поддерживает динамические операции вставки, удаления и перебалансировки с логарифмической сложностью, что позволяет работать с картами, содержащими миллионы точек, в реальном времени на ограниченных вычислительных ресурсах (например, на бортовом компьютере дрона). Это особенно важно при использовании лидаров Livox, чей неповторяющийся паттерн сканирования приводит к естественному накоплению плотной карты с течением времени.

\subsection{Livox как целевой сенсор}

Лидары Livox с неповторяющимся паттерном сканирования являются целевой аппаратной платформой для FAST-LIO. В отличие от механических лидаров с регулярной сеткой лучей, сканы Livox имеют нерегулярную пространственную структуру, что усложняет извлечение геометрических признаков (рёбер и плоскостей) традиционными методами, применяемыми в LOAM. FAST-LIO решает эту проблему путём прямой регистрации точек скана относительно карты (point-to-map registration) без явного выделения признаков, что делает алгоритм инвариантным к паттерну сканирования~\cite{xu2021,xu2022}.

\subsection{Ключевые параметры конфигурации}

Конфигурация FAST-LIO включает ряд параметров, влияющих на точность и производительность. Параметр \texttt{filter\_size\_surf} определяет размер воксельного фильтра для прореживания точек входного скана; \texttt{filter\_size\_map}~--- размер вокселя карты в ikd-Tree; \texttt{cube\_len} задаёт размер области поиска соответствий; \texttt{blind}~--- минимальную дистанцию от сенсора, точки ближе которой игнорируются (для фильтрации помех от элементов конструкции робота). Корректный подбор параметров критически важен: например, слишком большое значение \texttt{blind} приводит к потере полезной информации в indoor-сценах с близко расположенными стенами, а слишком маленькое~--- к включению в карту точек корпуса робота~\cite{xu2021}. Репозиторий алгоритма с открытым исходным кодом доступен на GitHub.\footnote{\url{https://github.com/hku-mars/FAST_LIO}} Исходный код написан на C++ с использованием библиотек Eigen~3 (линейная алгебра), PCL (Point Cloud Library, работа с облаками точек) и ROS (передача сообщений, визуализация)~\cite{xu2021}.

% ============================================================
\section{Датасеты HKU}

Авторами FAST-LIO опубликован набор открытых датасетов, записанных с использованием сенсорной платформы на базе лидара Livox AVIA (полупроводниковый лидар с неповторяющимся паттерном, дальность до 450~м при 80\,\% отражательной способности, FOV $70{,}4^\circ\times77{,}2^\circ$) и промышленного инерциального модуля BMI088 (частота опроса 200~Гц)~\cite{xu2021,xu2022}. Датасеты распространяются в формате rosbag и включают следующие сцены:
\begin{itemize}
  \item \texttt{quick-shack} (indoor): небольшое помещение ($\sim7\times7$~м) с характерной геометрией~--- стены, шкафы, препятствия. Короткая траектория длительностью около 30~с. Используется для быстрой верификации работоспособности алгоритма.
  \item \texttt{hku\_mb} (indoor): коридор главного корпуса HKU. Протяжённый indoor-маршрут с возвратно-поступательным движением вдоль длинного коридора. Характерен наличием повторяющихся структур (двери, окна) и стеклянных поверхностей, создающих сложности для лазерного сканирования.
  \item \texttt{mainbuilding 10Hz} (outdoor): периметр главного здания HKU. Запись с частотой сканирования 10~Гц. Типичная outdoor-сцена с фасадами зданий, элементами благоустройства, растительностью и дорожным покрытием.
  \item \texttt{mainbuilding 100Hz} (outdoor): та же траектория, что и \texttt{mainbuilding 10Hz}, но с максимальной частотой сканирования 100~Гц. Позволяет оценить влияние частоты сканирования на качество и вычислительную нагрузку одометрии.
\end{itemize}

Каждый rosbag-файл содержит следующие топики: \texttt{/livox/lidar} с типом сообщений \texttt{livox\_ros\_driver/CustomMsg} (облака точек в формате Livox) и \texttt{/livox/imu} с типом \texttt{sensor\_msgs/Imu} (данные инерциального модуля). Частота публикации облаков точек составляет 10~Гц для большинства датасетов и 100~Гц для высокочастотной версии \texttt{mainbuilding}. Датасеты не включают эталонной (ground truth) траектории, записанной от внешней системы трекинга; качество одометрии оценивается косвенно~--- по величине дрейфа между начальной и конечной точкой траектории, которая на замкнутых маршрутах теоретически равна нулю (метрика start-end position error).

% ============================================================
\section{Методика исследования влияния зеркалирования}

Для предварительного исследования влияния дополнительной плотности облака точек на точность лидарной одометрии проведён вычислительный эксперимент с синтетическим зеркалированием. Идея эксперимента состоит в имитации эффекта пассивного зеркала путём добавления к точкам исходного облака их отражённых копий с контролируемой долей (плотностью зеркалирования). Методика включает следующие этапы.

\subsection{Генерация модифицированных данных}

Для каждого из четырёх исходных датасетов создаются версии с различной долей зеркалированных точек: 5\,\%, 10\,\%, 20\,\%, 30\,\% и 50\,\%. Базовый уровень (0\,\%) соответствует исходным данным без модификаций.

\subsection{Математическая модель зеркалирования}

Для выбранной доли случайных точек $p_i=(x_i,y_i,z_i)$ применяется планарное отражение относительно виртуальной плоскости зеркала:
\begin{equation}
  y_i \rightarrow 2d - y_i,
  \label{eq:mirror}
\end{equation}
где $d=0{,}5$~м~--- расстояние от оптического центра сенсора до плоскости зеркала. Координаты $x_i$ и $z_i$ не изменяются. Данная модель соответствует оператору Хаусхолдера для частного случая вертикального зеркала с нормалью $n=(0,1,0)^T$~\cite{vishnevsky2026}:
\begin{equation}
  H = I - 2nn^T, \qquad H \in \mathbb{R}^{4\times4},
  \label{eq:householder}
\end{equation}
где $I$~--- единичная матрица, а в однородных координатах оператор дополняется членом $2dn$.

\subsection{Обработка и метрика качества}

Каждая модифицированная версия датасета обрабатывается алгоритмом FAST-LIO с идентичными параметрами конфигурации (единый launch-файл, одинаковые значения \texttt{filter\_size\_surf}, \texttt{filter\_size\_map}, \texttt{blind} и др.). Для каждого прогона вычисляется метрика start-end position error~--- евклидово расстояние между трёхмерными координатами начальной и конечной точек траектории, восстановленной из логарифма поз. На замкнутых маршрутах данная метрика характеризует накопленный дрейф одометрии. Всего выполнено 24 эксперимента (4 датасета $\times$ 6 уровней плотности, включая базовый 0\,\%).

% ============================================================
% \section{Эксперимент}
% Раздел будет добавлен после завершения серии экспериментов.
% Планируется: comparative evaluation of FAST-LIO accuracy on original
% and synthetically mirrored point clouds, tables and plots of
% start-end position error versus mirror density, discussion of
% the optimal density range and limitations.

% ============================================================
% \section{Заключение}
% Будет сформулировано после получения полных экспериментальных данных.

% ============================================================
% Список литературы — вариант 1: вручную
\begin{thebibliography}{9}

\bibitem{vishnevsky2026} Вишневский М.~Р., Фёдоров И.~Ю. Повышение пространственного разрешения механических LiDAR-систем с использованием виртуального зеркального канала и автоматической калибровки. \textit{Труды Института системного программирования РАН}, т.~38, №~3, с.~101--114, 2026. DOI: 10.15514/ISPRAS-2026-38(3)-38.

\bibitem{epifanov2024} Епифанов К.~А. Исследование алгоритмов замыкания цикла в лидарной одометрии. \textit{Сборник научных трудов}, с.~57, 2024.

\bibitem{lidar3d2025} Development of a LiDAR-based 3D Scanner Using a Plane Mirror Base for Improved Scanning Coverage and ICP-based Pose Graph SLAM Optimization. In \textit{Proceedings of ELTICOM}, 2025. DOI: 10.1109/ELTICOM67568.2025.11335854.

\bibitem{flatness2021} Laser Scanning Based Surface Flatness Measurement Using Flat Mirrors for Enhancing Scan Coverage Range. \textit{Remote Sensing}, vol.~13, no.~4, art.~714, 2021. DOI: 10.3390/rs13040714.

\bibitem{xu2021} Xu W., Zhang F. FAST-LIO: A Fast, Robust LiDAR-inertial Odometry Package by Tightly-Coupled Iterated Kalman Filter. \textit{IEEE Robotics and Automation Letters}, vol.~6, no.~2, pp.~3317--3324, 2021. DOI: 10.1109/LRA.2021.3064227.

\bibitem{xu2022} Xu W., Cai Y., He D., Lin J., Zhang F. FAST-LIO2: Fast Direct LiDAR-Inertial Odometry. \textit{IEEE Transactions on Robotics}, vol.~38, no.~4, pp.~2053--2073, 2022. DOI: 10.1109/TRO.2022.3141876.

\bibitem{thrun2005} Thrun S., Burgard W., Fox D. Probabilistic Robotics. MIT Press, 2005. 477~p.

\bibitem{zhang2014} Zhang J., Singh S. LOAM: Lidar Odometry and Mapping in Real-time. In \textit{Robotics: Science and Systems (RSS)}, 2014.

\bibitem{shan2018} Shan T., Englot B. LeGO-LOAM: Lightweight and Ground-Optimized Lidar Odometry and Mapping on Variable Terrain. In \textit{IEEE/RSJ IROS}, 2018, pp.~4758--4765.

\bibitem{shan2020} Shan T., Englot B., Meyers D., Wang W., Ratti C., Rus D. LIO-SAM: Tightly-coupled Lidar Inertial Odometry via Smoothing and Mapping. In \textit{IEEE/RSJ IROS}, 2020, pp.~5135--5142.

\bibitem{beraldin2010} Beraldin J.-A., Blais F., Lohr U. Laser Scanning Technology. In \textit{Airborne and Terrestrial Laser Scanning}, Vosselman G., Maas H.-G. (eds.). Whittles Publishing, 2010, pp.~1--42.

\bibitem{wang2020} Wang D., Watkins C., Xie H. MEMS Mirrors for LiDAR: A Review. \textit{Micromachines}, vol.~11, no.~5, art.~456, 2020. DOI: 10.3390/mi11050456.

\bibitem{yoo2018} Yoo H.~W., Druml N., Brunner D. et al. MEMS-based Lidar for Autonomous Driving. \textit{Elektrotechnik und Informationstechnik}, vol.~135, no.~6, pp.~408--415, 2018.

\bibitem{besl1992} Besl P.~J., McKay N.~D. A Method for Registration of 3-D Shapes. \textit{IEEE Transactions on Pattern Analysis and Machine Intelligence}, vol.~14, no.~2, pp.~239--256, 1992.

\bibitem{segal2009} Segal A., Hahnel D., Thrun S. Generalized-ICP. In \textit{Robotics: Science and Systems (RSS)}, 2009.

\end{thebibliography}

\end{document}
